% Part:sets-functions-relations
% Chapter: size-of-sets
% Section: non-enumerability-alt

\documentclass[../../../include/open-logic-section]{subfiles}

\begin{document}

\olfileid{sfr}{siz}{nen-alt}

\olsection{\printtoken{S}{nonenumerable} సమితులు}

\begin{editorial}
  ఈ విభాగం \olref[enm-alt]{sec}లోని నిర్వచనాలను వాడి
  $\Bin^\omega$, $\Pow{\Nat}$ల లెక్కించలేనితనాన్ని నిరూపిస్తుంది;
  అంటే, $\PosInt$ నుంచి సంగ్రస్త ప్రమేయానికి బదులు $\Nat$తో ద్విగుణ
  ప్రమేయాన్ని కోరుతుంది.
\end{editorial}

\begin{explain}
సహజ సంఖ్యల సమితి $\Nat$ అనంతమైనది. అది తేలికగానే !!{enumerable}
కూడా. అయితే \emph{!!{nonenumerable}} సమితులు, అంటే !!{enumerable}
కాని సమితులు ఉండటం విశేషమైన వాస్తవం
(\olref[sfr][siz][enm-alt]{defn:enumerable} చూడండి).

ఇది ఆశ్చర్యంగా అనిపించవచ్చు. ఎందుకంటే $A$ !!{nonenumerable} అని
చెప్పడం అంటే !!{bijection} $f \colon \Nat \to A$ \emph{ఏదీ లేదని}
చెప్పడమే; అంటే, $\Nat$లోని అనంతమైన !!{element}sను $A$లోకి పంపే ఏ
ప్రమేయమూ $A$ మొత్తాన్ని అయిపోగొట్టదు. కాబట్టి $A$ !!{nonenumerable}
అయితే, సహజ సంఖ్యల కంటే $A$లో ``ఎక్కువ'' !!{element}s ఉన్నాయి.

ఒక సమితి !!{nonenumerable} అని నిరూపించాలంటే, తగిన !!{bijection}
ఏదీ ఉండజాలదని చూపాలి. అందుకు ఉత్తమ మార్గం, $A$ యొక్క !!{element}sను
లెక్కించే ప్రతి ప్రయత్నమూ కనీసం ఒక !!{element}ను వదిలేస్తుందని
చూపడం; దీనివల్ల $f\colon \Nat \to A$ అనే ఏ ప్రమేయమూ !!{surjective}
కాదని తెలుస్తుంది. దీన్ని నిర్ధారించే ఒక సాధారణ వ్యూహం కాంటర్
\emph{వికర్ణ పద్ధతి}. $A$ యొక్క !!{element}s జాబితా $x_1$, $x_2$, \dots{}
ఇచ్చినప్పుడు, దాని నిర్మాణం వల్లనే ఆ జాబితాలో ఉండజాలని $A$ యొక్క
మరో !!{element}ను నిర్మిస్తాం.

ఇదంతా ఉదాహరణతోనే బాగా అర్థమవుతుంది. కాబట్టి మన మొదటి ఉదాహరణ $0$,
$1$ల అనంత పదబంధాలన్నిటి సమితి~$\Bin^\omega$. (`$\Bin$' అంటే ద్విమితి;
దాన్ని రెండు-మూలకాల సమితి $\{0,1\}$గా భావించవచ్చు.)
\oliflabeldef{sfr:card-arithmetic:card-opps:sec}{\footnote{మరింత
కచ్చితంగా, $\Bin^\omega$ను $0$, $1$ల $\omega$-క్రమాలన్నిటి సమితి,
అంటే $\funfromto{\omega}{\{0,1\}}$ అని నిర్దేశించాలి. కానీ దీని అర్థం
\olref[sfr][card-arithmetic][card-opps]{sec}లో మాత్రమే స్పష్టమవుతుంది.}}{}
అయితే ప్రస్తుతానికి ఈ కొంచెం సడలించిన వివరణ వల్ల గందరగోళం
రాకూడదు.
\end{explain}

\begin{thm}
\ollabel{thm:nonenum-bin-omega}
$\Bin^\omega$~!!{nonenumerable}.
\end{thm}

\begin{proof}
$\Bin^\omega$ యొక్క ఒక ఉపసమితికి ఏదైనా లెక్కింపును పరిశీలించండి.
అంటే, $s_{0}$, $s_{1}$, $s_{2}$, \dots{} అనే జాబితా ఉంది; ఇందులో
ప్రతి $s_n$, $0$, $1$ల అనంత పదబంధం. $s_n(m)$ అంటే ఈ జాబితాలోని
$n$వ పదబంధంలోని $m$వ అంకె అనుకుందాం. అప్పుడు మన జాబితిని ఒక
పట్టికగా భావించవచ్చు; అందులో $s_n(m)$ను $n$వ అడ్డువరుసలో, $m$వ
నిలువువరుసలో పెడతాం:
\sourcecorrection{OLSIZ-008}{ఆంగ్ల మూలం $s_n(m)$ను $m$వ పదబంధంలోని
$n$వ అంకెగా చెప్పి, వెంటనే దాన్ని $n$వ అడ్డువరుస, $m$వ నిలువువరుసలో
ఉంచుతుంది. పట్టిక, సంకేతాలకు సరిపడేలా ఇక్కడ $n$వ పదబంధంలోని $m$వ
అంకె అని సరిచేశాం.}
\[
\begin{array}{c|c|c|c|c|c}
& 0 & 1 & 2 & 3 & \dots \\\hline
0 & \mathbf{s_{0}(0)} & s_{0}(1) & s_{0}(2) & s_0(3) & \dots \\\hline
1 & s_{1}(0)& \mathbf{s_{1}(1)} & s_1(2) & s_1(3) & \dots \\\hline
2 & s_{2}(0)& s_{2}(1) & \mathbf{s_2(2)} & s_2(3) & \dots \\\hline
3 & s_{3}(0)& s_{3}(1) & s_3(2) & \mathbf{s_3(3)} & \dots \\\hline
\vdots & \vdots & \vdots & \vdots & \vdots & \mathbf{\ddots}
\end{array}
\]
ఇప్పుడు ఈ జాబితాలో లేని $0$, $1$ల అనంత పదబంధం $d$ను నిర్మిస్తాం.
దాని ప్రతి నమోదును పేర్కొనడం ద్వారా, అంటే ప్రతి $n \in \Nat$కు
$d(n)$ను పేర్కొనడం ద్వారా దీన్ని చేస్తాం. సహజంగా చెప్పాలంటే, పై
పట్టిక వికర్ణం వెంబడి చదివి (అందుకే ``వికర్ణ పద్ధతి'' అనే పేరు),
ప్రతి $1$ను $0$గానూ, ప్రతి $0$ను $1$గానూ మారుస్తాం.
\sourcecorrection{OLSIZ-009}{ఆంగ్ల మూలం రెండు సందర్భాల్లోనూ $1$ను
$0$గా మార్చమని పునరావృతం చేస్తుంది. వెంటనే వచ్చే సందర్భాలవారీ
నిర్వచనానికి సరిపడేలా రెండవ సందర్భాన్ని $0$ నుంచి $1$కు మార్చాం.}
మరింత అమూర్తంగా, వికర్ణంలోని $n$వ !!{element} $s_n(n)$ విలువ $1$
లేదా $0$ అవడాన్ని బట్టి $d(n)$ను $0$ లేదా $1$గా నిర్వచిస్తాం:
\[
d(n) =
\begin{cases}
1 & \text{$s_{n}(n) = 0$ అయితే}\\
0 & \text{$s_{n}(n) = 1$ అయితే}
\end{cases}
\]
ఇది $0$, $1$ల అనంత పదబంధం కనుక స్పష్టంగా $d \in \Bin^\omega$.
కానీ ప్రతి $n \in \Nat$కూ $d(n) \neq s_n(n)$ అయ్యేలా $d$ను
నిర్మించాం. అంటే $d$, $s_n$ నుంచి దాని $n$వ నమోదులో భిన్నంగా ఉంటుంది.
కాబట్టి ప్రతి $n\in \Nat$కూ $d \neq s_n$. అందువల్ల $d$, $s_0$,
$s_1$, $s_2$, \dots జాబితాలో ఉండజాలదు.

$\Bin^\omega$ యొక్క ఏదో ఉపసమితి యొక్క యాదృచ్ఛిక లెక్కింపు ఇచ్చినా,
అది $\Bin^\omega$ యొక్క ఏదో !!{element}ను వదిలేస్తుందని చూపాం.
కాబట్టి $\Bin^\omega$ సమితికి లెక్కింపు లేదు; అంటే $\Bin^\omega$
!!{nonenumerable}.
\end{proof}

\begin{explain}
ఈ నిరూపణ పద్ధతి $d$ను నిర్వచించడానికి పట్టిక వికర్ణాన్ని వాడుతుంది
కనుక దీన్ని ``వికర్ణీకరణ'' అంటారు. అయితే వికర్ణీకరణకు పట్టిక ఉండాల్సిన
అవసరం లేదు. నిజానికి, పట్టిక గానీ అసలు వికర్ణం గానీ లేకపోయినా ఇలాంటి
ఆలోచనతో ఏదో సమితి !!{nonenumerable} అని చూపవచ్చు. కింది ఫలితం దాన్ని
వివరిస్తుంది.
\end{explain}

\begin{thm}
\ollabel{thm:nonenum-pownat}
$\Pow{\Nat}$ !!{enumerable} కాదు.
\end{thm}

\begin{proof}
$\Nat$ ఉపసమితుల ప్రతి జాబితా $\Nat$ యొక్క ఏదో ఉపసమితిని వదిలేస్తుందని
చూపుతూ ఇదే విధంగా ముందుకు వెళ్తాం. $\Nat$ ఉపసమితుల జాబితా
$N_0, N_1, N_2, \ldots$ ఉందనుకుందాం. $n \in D$ అవడానికి అవసరమైన
మరియు సరిపడే షరతు $n \notin N_{n}$ అయ్యేలా $D$ను నిర్వచిద్దాం:
\[
D = \Setabs{n \in \Nat}{n \notin N_n}
\]
స్పష్టంగా $D\subseteq \Nat$. కానీ $D$ జాబితాలో ఉండజాలదు. ఎందుకంటే
నిర్మాణం ప్రకారం $n \in D$ అవడానికి అవసరమైన మరియు సరిపడే షరతు
$n\notin N_n$ అవడం; కాబట్టి ప్రతి $n \in \Nat$కూ $D \neq N_n$.
\end{proof}

\begin{explain}
ముందటి నిరూపణ వికర్ణాన్ని ప్రస్తావించలేదు. అయినా దాన్ని ఇలా చిత్రీకరిస్తే,
అందులో వికర్ణం ఉందని భావించవచ్చు: $N_0$, $N_1$, \dots సమితులను ఒక
పట్టికలో రాసినట్లు ఊహించండి; $m \in N_n$ అయితేనే $m$ను $m$వ
నిలువువరుసలో రాస్తూ, $N_n$ను $n$వ అడ్డువరుసలో రాయండి. ఉదాహరణకు,
ఆ జాబితాలోని మొదటి నాలుగు సమితులు $\{0,1,2,\dots\}$,
$\{1, 3, 5, \dots\}$, $\{0,1,4\}$, $\{2,3,4,\dots\}$ అనుకుందాం;
అప్పుడు మన పట్టిక ఇలా మొదలవుతుంది:
\[
\begin{array}{r@{}rrrrrrr}
  N_0 = \{ & \mathbf{0}, & 1, & 2, & & & & \dots\}\\
  N_1 = \{ &  & \mathbf{1}, &  & 3, &  & 5, & \dots\}\\
  N_2 = \{ & 0, & 1, &  &  & 4\phantom{,} &  & \}\\
  N_3 = \{ &  &  & 2, & \mathbf{3}, & 4, & & \dots\}\\
  &\vdots & & & & & & \ddots\phantom{\}}
  \end{array}
\]
వికర్ణం వెంబడి వెళ్లి, $n$ వికర్ణంపై \emph{లేకపోతేనే}
$n \in D$గా ఉంచి పొందిన సమితి $D$. కాబట్టి పై సందర్భంలో $0$, $1$లను
వదిలేస్తాం, $2$ను చేర్చుతాం, $3$ను వదిలేస్తాం, తదితరంగా.
\end{explain}

\begin{prob}
$f \colon \Nat \to \Nat$ రూపంలోని ప్రమేయాలన్నిటి సమితి
!!{nonenumerable} అని ప్రత్యక్ష వికర్ణ వాదనతో చూపండి. అంటే,
$f_1$, $f_2$, \dots{} ప్రమేయాల జాబితి, ప్రతి $f_i\colon
\Nat \to \Nat$ అయినప్పుడు, ఈ జాబితాలో లేని ఏదో
$g \colon \Nat \to \Nat$ ఉందని చూపండి.
\end{prob}

\end{document}
